\bookchapter{第十一章：撮影と画像管理}

\chapterintro{一、あなたもこれらの間違いをしていませんか？}

\booksection{まず最も一般的な問題をいくつか挙げます。いくつ当てはまるか確認してください：}

レンズが汚れている

\begin{itemize}
  \item \textbf{：これが本当に最も一般的な問題。スマホを取り出す時に指紋がレンズにつき、撮影結果が霞んで、スマホが壊れたと思い込みがち。習慣づけて、撮影前に服の裾でレンズを拭きましょう。}露出を調整しない
  \item \textbf{：デフォルトモードでは、画面をタップするとフォーカスされますが、タップ後に上下にスワイプすると明るさを調整できます。多くの人が撮った写真が暗すぎたり明るすぎたりするのは、この操作を知らないからです。}フラッシュを常にオン
  \item \textbf{：スマホのLEDフラッシュで撮った効果は本当に見苦しく、光が硬くて冷たく、人が幽霊のように写ります。本当に真っ暗闇でない限り、オフにしましょう。暗所では、夜景モードの方がフラッシュよりはるかに効果的です。}ピンチアウトで拡大
  \item \textbf{：画面上で二本指ピンチはデジタルズームで、本質的に画像を切り抜いて拡大するだけなので、画質が明確に低下します。遠くを撮りたい場合は、ファインダー上のズームボタン（0.5x/1x/2x/3x）を直接使いましょう。これらは光学ズームで画質を損ないません。}：画面上でピンチズームするとデジタルズームになり、本質的には画像を切り抜いて拡大するため画質が明らかに低下します。遠くを撮りたい場合は、ファインダー上のズームボタン（0.5x/1x/2x/3x）を直接使ってください。0.5x（超広角）と1x（メインカメラ）はネイティブ光学レンズで画質劣化がありません。2x/3xは「光学品質ズーム」——高画素センサーの切り抜きで実現し、画質は光学ズームに近いですが独立レンズではありません。3xを超えるズームは通常ハイブリッドズームまたはデジタルズームで、画質が明らかに低下します。
\end{itemize}

\booksection{これらは理論知識不要、一度使えば効果がわかります：}

三分割法構図＋グリッド線

\begin{itemize}
  \item \textbf{：カメラ設定で「グリッド線」（ガイド線とも）をオン。撮影時に強調したいものを4本の線の交点に置くと、画面がすぐに「正しく」なります。これは写真で最もシンプルかつ最も効果的なルールで、神秘性ではなく、人間の視覚習慣に合致するだけです。}「リーディングライン」を探す
  \item \textbf{：風景や街並みを撮る時、画面内に自然な線がないか留意——道路、手すり、壁、並木——これらの線を画面の端から被写体まで伸ばします。このテクニックを使いこなすと、通りすがりの人があなたの写真を見て写真を学んだことがあると思います。}立ったまま撮らない
  \item \textbf{：多くの人が撮影する時は常に目の高さなので、視点がワンパターン。しゃがんで低アングルから撮ったり、高いところから俯瞰撮影してみたりすると、同じ場所でも全く違う印象に。子供やペットを撮る時は特にしゃがむ必要があり、目線の高さで撮ると自然な感じになります。}光が第一要素
  \item \textbf{：真昼の太陽の下でポートレートを撮ると、顔の影が硬くて醜くなります。最も良い光は早朝と夕方の「ゴールデンアワー」で、太陽が斜めで色が暖色寄り、何を撮っても美しくなります。曇りの日の光は地味ですが、均一で柔らかく、ポートレートは真昼の太陽より撮りやすいです。}夜景モードの使い方
  \item \textbf{：夜景モードをオンにすると、スマホが自動的に複数枚連写して合成します。この時に最も重要なのは——絶対に手ブレしないこと。壁に寄りかかれるなら寄りかかり、机に置けるなら置き、ダメなら息を止めましょう。夜景モードは2〜4秒安定を保つ必要があり、ブレると画面がぼやけます。}三、ポートレートモードとプロモード
\end{itemize}

\booksection{ポートレートモード}

\begin{itemize}
  \item \textbf{：背景がぼけて被写体がくっきりする効果（浅い被写界深度）を撮れます。スマホはレンズの物理サイズが小さいため、本物のボケ能力は限定的で、AIアルゴリズムで実現しています——簡単に言うと、AIが人物を切り抜き、背景をぼかします。使用時は1〜2メートルの距離が最も効果的で、遠すぎても近すぎてもアルゴリズムがエラーを起こしやすい。またポートレート撮影時は、背景に複雑な線がないよう注意。そうでないとAIの切り抜きで髪の辺りに「失敗」が生じます。}プロモード
  \item \textbf{：画面効果を完全にコントロールしたい場合は、プロモードを試しましょう。いくつかの重要なパラメータを知っておく必要があります：ISOが低いほど画質がクリーンですが、画面は暗くなります。シャッターが遅いほど光量が増えますが、手ブレでぼやけます。ホワイトバランスは色の暖かさ冷たさを管理し、調整が悪いと写真が青や黄色に偏ります。一般人は常時プロモードを使う必要はありませんが、これらのパラメータを理解すれば、カメラが何をしているか把握できます。}四、レタッチ：やりすぎず、ちょうど良く
\end{itemize}

\booksection{多くの人がレタッチについて誤解しており、レタッチすると不自然になると思っています。実はカメラで撮った写真は元々肉眼で見たものと異なり、適度な調整は当時見た画面を再現するためで、偽造ではありません。}

Snapseed

\begin{itemize}
  \item \textbf{：Google製の無料レタッチツール、機能はプロ級だが広告なし。私の日常的なレタッチフローは：まず「画像調整」で明るさとコントラストを調整 「ディテール」でシャープ化 「ホワイトバランス」で色かぶり補正 最後に「切り抜き」で希望の構図に。Snapseedの「局部」機能は便利で、画面内の特定領域だけを調整でき、全画像を一括変更する必要がありません。}Xingtu / Meitu
  \item \textbf{：この2つは人物写真のレタッチに適しています。ただし原則として：美顔フィルターは30\%以内に抑えるべきで、それを超えると不自然になり、肌がプラスチックのようになります。完璧より自然な方が美しい。}Lightroomモバイル版
  \item \textbf{：色にこだわるなら、これが最もプロフェッショナル。トーンカーブ調整、HSL色別調整などの高度な機能があります。ただし多くの人には大げさすぎるかもしれません。}レタッチワークフローの提案
  \item \textbf{：私の習慣は、まず良い原画像を撮り、帰宅後にパソコン画面やiPadでじっくりレタッチすること。スマホ画面は小さく、色もあまり正確ではありません。どうしてもスマホでレタッチする必要がある場合は、画面の明るさを中〜上程度に調整し、省電力モードでレタッチしないように。}五、写真が多すぎる？整理戦略
\end{itemize}

\booksection{スマホのアルバムに数千枚の写真が溜まり、特定の一枚を探すのに時間がかかる——この経験は皆あるでしょう。いくつかの実用的な整理方法：}

検索を活用

\begin{itemize}
  \item \textbf{：最近のiPhoneと多くのAndroidスマホのアルバムはスマート検索に対応。「猫」「海辺」「身分証明書」「花」などのキーワードで検索すると、システムがAIで画像内容を認識して探し出します。日付や場所でも検索でき、例えば「昨年5月」「北京」。この機能を知る人は少ないですが、本当に使いやすい。}定期的にゴミを掃除
  \item \textbf{：スクショ、重複写真、ピンボケの失敗作、これらが容量を食う元凶。iOSアルバムには「重複項目」自動統合機能があり、Androidはファイルマネージャーで「screenshots」フォルダを検索し、不要なスクショを一括削除できます。毎月5分かけて掃除するだけで数GB空きます。}アルバムを作成、デフォルト並び順だけに頼らない
  \item \textbf{：システムデフォルトの時系列並び順は間違いではないですが、「子供の成長写真」「旅行写真」「重要書類」のような時期をまたぐコンテンツは、手動でアルバムを作って分類すべき。Appleの「共有アルバム」は家族と一緒に写真を追加でき、一人が作成すれば家族全員が閲覧できます。}バックアップは最低限の要件
  \item \textbf{：写真は最も代替不能なデータ——ファイルは再構築できますが、チャット履歴はなくなればそれまで、でも子供の写真や老人の動画はなくなると本当に戻りません。必ずクラウドバックアップを有効に。iCloud写真が最も手軽、Androidは各社独自のクラウドサービスがあるか、Baiduネットディスク/Alibabaクラウドドライブで自動バックアップ。原則は：少なくとも1つのクラウドバックアップを持ち、卵を一つのカゴに盛らないこと。}六、スクショと画面録画
\end{itemize}

\booksection{スクショ}

\begin{itemize}
  \item \textbf{：電源キー＋音量上げキーの同時押し、すべてのスマホで共通の方法。多くのAndroidスマホは三指下スワイプスクショ、関節ダブルタップスクショにも対応。長尺スクショ（スクロールスクショ）はAndroidで成熟——1枚スクショ後に「スクロール」ボタンをタップすると、スマホが自動的にページを下に送って長尺画像に結合。チャット履歴やウェブページのスクショに適しています。}画面録画
  \item \textbf{：コントロールセンターに「画面録画」を追加し、タップで録画開始、もう一度タップで停止。録画前に長押しするとマイク音声を録画するか選択可能——チュートリアル解説ならマイクオン、操作画面だけならオフ。録画完了した動画は自動的にアルバムに保存されます。}練習：スマホで撮影＋レタッチ＋整理の完全フローを一度実行
\end{itemize}

\begin{tcbpractice}{カメラ設定で「グリッド線」をオンにし、自然光のある場所で3枚撮影（レンズ拭き、露出調整に注意）。}
\begin{enumerate}
  \item 最も良い1枚を選び、Snapseedで基本調整（明るさ、コントラスト、切り抜き）。
  \item この写真に専用アルバムを作成し、「私の作品」と命名。
  \item クラウドバックアップが有効か確認し、写真がクラウドに1部保存されていることを確保。
  \item 第十一章付録：プロフェッショナル映像技術知識
\end{enumerate}

\end{tcbpractice}
